% USEFUL LINKS:
% -------------
%
% - UiO LaTeX guides:          https://www.mn.uio.no/ifi/tjenester/it/hjelp/latex/
% - Mathematics:               https://en.wikibooks.org/wiki/LaTeX/Mathematics
% - Physics:                   https://ctan.uib.no/macros/latex/contrib/physics/physics.pdf
% - Basics of Tikz:            https://en.wikibooks.org/wiki/LaTeX/PGF/Tikz
% - All the colors!            https://en.wikibooks.org/wiki/LaTeX/Colors
% - How to make tables:        https://en.wikibooks.org/wiki/LaTeX/Tables
% - Code listing styles:       https://en.wikibooks.org/wiki/LaTeX/Source_Code_Listings
% - \includegraphics           https://en.wikibooks.org/wiki/LaTeX/Importing_Graphics
% - Learn more about figures:  https://en.wikibooks.org/wiki/LaTeX/Floats,_Figures_and_Captions
% - Automagic bibliography:    https://en.wikibooks.org/wiki/LaTeX/Bibliography_Management  (this one is kinda difficult the first time)
%
%                              (This document is of class "revtex4-1", the REVTeX Guide explains how the class works)
%   REVTeX Guide:              http://www.physics.csbsju.edu/370/papers/Journal_Style_Manuals/auguide4-1.pdf
%
%
% COMPILING THE .pdf FILE IN THE LINUX TERMINAL
% ---------------------------------------------
%
% [terminal]$ pdflatex report_example.tex
%
% Run the command twice, always.
%
% When using references, footnotes, etc. you should run the following chain of commands:
%
% [terminal]$ pdflatex report_example.tex
% [terminal]$ bibtex report_example
% [terminal]$ pdflatex report_example.tex
% [terminal]$ pdflatex report_example.tex
%
% This series of commands can of course be gathered into a single-line command:
% [terminal]$ pdflatex report_example.tex && bibtex report_example.aux && pdflatex report_example.tex && pdflatex report_example.tex
%
% ----------------------------------------------------



% \documentclass[english,notitlepage,reprint,nofootinbib]{revtex4-2}  % defines the basic parameters of the document
\documentclass[english,notitlepage,reprint,nofootinbib]{revtex4-2}  % defines the basic parameters of the document
% If you want a single-column, remove "reprint"

% Allows special characters (including æøå)
\usepackage[utf8]{inputenc}
% \usepackage[english]{babel}

% Note that you may need to download some of these packages manually, it depends on your setup.
% It may be usefult to download TeXMaker, because it includes a large library of the most common packages.

\usepackage{physics,amssymb}  % mathematical symbols (physics imports amsmath)
\include{amsmath}
\usepackage{graphicx}         % include graphics such as plots
\usepackage{xcolor}           % set colors
\usepackage{hyperref}         % automagic cross-referencing
\usepackage{listings}         % display code
\usepackage{subfigure}        % imports a lot of cool and useful figure commands
% \usepackage{float}
%\usepackage[section]{placeins}
\usepackage{algorithm}
\usepackage[noend]{algpseudocode}
\usepackage{subfigure}
\usepackage{tikz}
\usetikzlibrary{quantikz}
% defines the color of hyperref objects
% Blending two colors:  blue!80!black  =  80% blue and 20% black
\hypersetup{ % this is just my personal choice, feel free to change things
    colorlinks,
    linkcolor={red!50!black},
    citecolor={blue!50!black},
    urlcolor={blue!80!black}}


% ===================== Shortcuts/new commands ======================
\renewcommand{\vec}[1]{\mathbf{#1}}


\begin{document}

\title{FYS4150 - Project 3\\[0.5em]
  \small \href{https://github.uio.no/icrukan/FYS4150/tree/master/Project3}{\texttt{github.uio.no/icrukan/FYS4150/Project3}}
}
\author{Magnus Axelsen and Isak Cecil Onsager Rukan} % self-explanatory
\date{\today}                             % self-explanatory
\noaffiliation                            % ignore this, but keep it.
\maketitle

\section{Introduction}
... 

\section{Theory}
\subsection{The Penning Trap}
A penning trap is a device which `traps' an electrically charged particle inside a static (time-independent) electromagnetic field. How this is achieved can be understood through the Lorentz force,
\begin{align}
    \vec{F} = q\vec{E} + q\vec{v}\times \vec{B},    \label{eq:Lorentz_force}
\end{align}
which describes the force experienced by a particle of charge \(q\) inside an electromagnetic field. Here \(\vec{v}\) is the velocity of the particle, \(\vec{E}\) the electric field, and \(\vec{B}\) the magnetic field. For a positively charged particle, the Lorentz force will apply a force in the direction of the \(E\)-field, as well as a force perpendicular to the particle's velocity and the direction of the \(B\)-field. The idea of Penning trap is then to make a particle bounded, say vertically in the \(z\)-plane, by a particular setup of the \(E\)-field, and rotating in the \(xy\)-plane through the second term in \eqref{eq:Lorentz_force}.

A Penning trap can be realized by a static and homogeneous \(B\)-field pointing in the \(z\)-direction, together with two positively charged electrodes placed above a negatively charged electrode, with the latter having a circular shape in the \(xy\)-plane. The setup is described in Fig.~\ref{fig:penning_trap} as a cross section in the \(zx\)- (or \(zy\)-)plane.
\begin{figure}[h!]
    \centering
    \includegraphics[scale=0.09]{figs/Penning_Trap.pdf}
    \caption{A positively charged particle is trapped within three electrodes, creating a Penning trap. The negatively charged electrode (b) has a circular-like shape, and the illustration shows a cross section of (b). The two electrodes placed above each other, both labelled (a), are positively charged, and an external homogeneous \(B\)-field is applied in the vertical direction. The illustration is created by Arian Kriesch and can be found in \href{https://en.wikipedia.org/wiki/File:Penning_Trap.svg}{Wikimedia Commons}.}
    \label{fig:penning_trap}
\end{figure}  
Each electrode is held at an electric potential \(V_0\), and the distance between them \(d\), is referred to as the \emph{characteristic length}. The potential created by such electrodes can be described as 
\begin{align}
    V(x,y,z) = \frac{V_0}{2d^2}\left(2z^2 - x^2 -y^2\right),
\end{align}
and so the \(E\)-field becomes
\begin{align}
    \vec{E} = -\nabla V = \frac{V_0}{d^2}\left(x, y, -2z\right).
\end{align}
For the homogeneous \(B\)-field, we write 
\begin{align}
    \vec{B} = B_0\hat{e}_z = (0,0,B_0),
\end{align}
where \(B_0>0\) is the strength of the field. 






\onecolumngrid
% \bibliographystyle{apalike}
\bibliographystyle{unsrt}
\bibliography{ref}


\end{document}